\label{sec:background}

\subsection{ACS 5-Year Estimates: Methodology and Reliability}

The American Community Survey (ACS) is a continuous demographic survey conducted
by the U.S.\ Census Bureau, collecting information from approximately 3.5~million
housing units annually~\citep{acs2020}.
The 5-year estimates aggregate five calendar years of data to produce reliable
tract-level estimates; single-year ACS estimates are not published at the census
tract level because the sample is too small.
The 2016--2020 five-year estimates (``2020 ACS 5-year'') represent the most
recent tract-level housing data available as of this writing and correspond
to the 2020 census tract boundaries used by \textsc{bisect}'s redistricting pipeline.

\paragraph{Margin of error.}
Because ACS estimates are survey-based, each published figure carries a margin
of error (MOE) at the 90\% confidence interval.
For housing-stock variables at the census tract level, MOEs are typically within
10--15\% of the estimate for tracts with population above 2,000.
Smaller tracts and tracts in rural areas may have larger relative MOEs.
For the purposes of the housing character vector, we treat the published point
estimates as the true values and note that MOE-based weighting of the similarity
metric is a potential extension.
All empirical results in this paper are designated Phase~2 pending ACS pipeline
extension and accordingly carry the $^\dagger$ dagger notation.

\paragraph{Tables used.}
Table~\ref{tab:acs-tables} lists the three ACS tables from which the housing
character components are derived.
All three are available at census tract geography for every state, published
in the ACS 5-year series.

\begin{table}[ht]
\centering
\caption{ACS 5-year tables used for housing character vector construction.}
\label{tab:acs-tables}
\small
\begin{tabularx}{\textwidth}{llX}
\toprule
Table & Name & Housing character component \\
\midrule
B25024 & Units in Structure & Single-family fraction, multifamily fraction \\
B25003 & Tenure & Owner-occupancy rate \\
B25035 & Median Year Structure Built & Housing vintage score \\
\bottomrule
\end{tabularx}
\end{table}

\subsection{Why Housing Type Predicts Community Structure}

\subsubsection{Homeownership and Civic Participation}

The relationship between homeownership and civic engagement is one of the most
studied questions in American political sociology.
DiPasquale and Glaeser~\citeyearpar{dipasquale1999} find in a national panel
that homeowners are significantly more likely to know their congressional
representative, participate in local organizations, and vote in local elections
than otherwise-identical renters.
The authors attribute the effect to the financial stake homeowners hold in
local government outcomes: school quality, property tax rates, zoning variances,
and infrastructure maintenance all directly affect property values in a way
they do not affect rental costs.

Fischel~\citeyearpar{fischel2001} elaborates this mechanism through the concept
of the ``homevoter hypothesis'': homeowners, unlike renters, cannot diversify
their housing investment, so they vote as if their home equity were their
primary financial asset.
They vote for school board candidates, against upzoning proposals, and for
infrastructure maintenance bonds with a consistency that non-homeowners do not
exhibit.
The political implication for redistricting is direct: a tract that is
predominantly owner-occupied has a recognizably different political interest
profile from an adjacent tract that is predominantly rental.
They are not members of the same community of interest on housing and land-use
policy.

\subsubsection{Neighborhood Stability and the Filtering Theory}

The lifecycle of the housing stock---its vintage---is another dimension of
community structure.
The ``filtering'' theory of housing markets~\citep{rosenthal2014} holds that
new construction at the high end of the market gradually filters down to lower
income occupants as the housing ages.
A tract built before 1940 and still predominantly owner-occupied is a community
that has resisted the filtering dynamic: it has maintained its character across
multiple generations of residents.
Historic preservation districts, which are legally recognized institutional
forms of this resistance, are typically composed of pre-1940 housing stock.

The community stability of pre-1960 neighborhoods is not merely a function
of the age of their buildings; it is a function of the long tenures of their
residents~\citep{weaver2020}.
Long-tenured homeowners have different representation interests from recent
arrivals: they are invested in the long-run trajectory of local institutions,
they have organized around neighborhood associations with multi-decade histories,
and they have fought and won (or lost) land-use battles that shaped the
current built environment.
Splitting such a neighborhood across congressional district boundaries fragments
a community of interest in the most literal sense.

\subsubsection{NIMBYism as a Community Signal}

Einstein, Palmer, and Glick~\citeyearpar{einstein2019} analyze 2015--2017
local planning board meeting minutes from Massachusetts and find that
participants who speak against housing development proposals are disproportionately
homeowners, older, and white, while the broader community they claim to represent
is more diverse and more renting.
Their finding---that homeowner communities systematically mobilize against
density increases---documents the political homogeneity of single-family
homeowner communities.
This homogeneity is precisely what the housing character vector is designed
to measure: tracts where the homeowner fraction and single-family fraction
are both high share a political interest profile that low-SF, high-rental
tracts do not.

The NIMBYism signal is partisan-neutral in origin.
Single-family homeowner communities with strong opposition to upzoning
are documented in liberal San Francisco~\citep{einstein2019},
moderate suburban New Jersey~\citep{fischel2001}, and conservative
Houston~\citep{nolo2022zoning}.
The housing character signal does not predict partisan vote share;
it predicts community cohesion on the housing-and-land-use dimension
of political interest.

\subsubsection{Implementation Position in the M-Track Framework}

The M-track framework (M.0) defines a family of community character signals
from M.1 (economic character) to M.7 (transit accessibility).
Table~\ref{tab:mtrack-position} places M.3 within that family.
M.3 occupies the residential built environment dimension, which is not covered
by any other M-track signal.
M.1 covers the workplace-side economic character; M.3 covers the residential-side
physical character.
M.6 covers administrative zone co-membership, which overlaps with M.3 in the
school district dimension (school district membership correlates strongly with
single-family homeowner neighborhoods) but is legally and computationally
distinct.

\begin{table}[ht]
\centering
\caption{Position of M.3 in the M-track signal taxonomy.}
\label{tab:mtrack-position}
\small
\begin{tabular}{lll}
\toprule
Paper & Signal & Data source \\
\midrule
M.1 & Economic character (workplace) & LODES WAC \\
M.2 & Land use character & NLCD \\
M.3 & Housing character (residential) & ACS B25024/B25003/B25035 \\
M.4 & Commuting shed & LODES OD \\
M.5 & Topographic character & USGS 3DEP \\
M.6 & Administrative zone co-membership & TIGER/Line special districts \\
M.7 & Transit accessibility & NTD / GTFS \\
\bottomrule
\end{tabular}
\end{table}
